\documentclass[book.tex]{subfiles}

\begin{document}
\chapter{Normalizing Flows}
\label{chap:normalizing_flows}
\noindent\rule{11cm}{0.4pt} \\
\textbf{Prerequisites:}  \autoref{chap:ml_basics}, \autoref{chap:probability}  \\
\textbf{Difficulty Level:} ** \\
\noindent\rule{11cm}{0.4pt}
\vspace{5mm}

%So far, we've primarily discussed machine learning models suitable for property prediction. That is, given some input, the models we've described construct a representation of the underlying structure which they then use to learn structure-property relationships from training data. This learned model is then used to make predictions of the properties of new inputs which haven't been seen in the model's training data.

%This framework is very powerful but has some limitations. For one, often the challenge we face as scientists is that we'd like to design a new molecule, or material or system which has some desired properties. How can we use machine learning to aid in this process of design?

%One simple strategy is to perform a combinatorial search. The basic idea here is that we construct an algorithm which enumerates many potentially interesting inputs. These hypothetical inputs can be fed into a property prediction deep network which will give us an estimate of the properties of this material.

%This combinatorial approach has a number of advantages. For one, it's simple and a good scientist can often construct a script that walks over a space of ideas. However, this model doesn't have scope for "creativity." In an ideal world, we'd like systems that serve as creativity amplifiers, helping us augment the space of ideas that we can explore. Generative deep models hold out the promise of helping us augment our creative facilities by enabling non-obvious algorithmically generated materials for us to test.

We'll start off this chapter with a mathematical review of the underpinnings of generative modeling. %In this process, we'll consider a number of different models such as variational modeling, generative adversarial networks, and normalizing flows. We'll see how we can use these powerful frameworks to help solve our motivating problem of materials design, but interestingly we'll see that generative modeling can be useful for other problems such as free energy estimation which are also of fundamental interest.

\section{Mathematical Underpinnings}

Let's consider a probability distribution $\mathcal{P}$ over $\mathbb{R}^N$. For now, we can informally represent this probability distribution by a probability density function $p: \mathbb{R}^n \to \mathbb{R}$. We previously introduced datasets by noting that training samples $x_i$ are sampled from this underlying data distribution  $\mathcal{P}$, which we write mathematically as $x_i \sim \mathcal{P}$. In some simple toy cases, $\mathcal{P}$ may be a tractable distribution which can be described analytically. However, even for analytical distributions, it can be quite challenging to draw samples from $\mathcal{P}$. Let's illustrate this with an example.

\subsection{Multivariate Gaussians}
The multivariate Gaussian distribution is a fundamental mainstay of the modern sciences. We typically write the distribution as $\mathcal{N}(\mu, \Sigma)$  where $\mu \in \mathbb{R}^N$ is the mean vector and $\Sigma \in \mathbb{R}^{N \times N}$ is the covariance matrix. Roughly, the multivariate Gaussian encodes a probability distribution centered around the mean $\mu$, where vectors that are closer to $\mu$ are more likely and where vectors further away grow progressively less likely. This rough intuition is formally encoded by the probability density function

\[ p_N(x) = \frac{1}{\sqrt{(2\pi)^N |\Sigma|}} \exp\left (-\frac{1}{2} (x - \mu)^T\Sigma^{-1}(x-\mu)\right ) \]

Here, $|\Sigma|$ is the determinant of the covariance matrix $\Sigma$. How can we sample from this  distribution? Creating a suitable sampling algorithm turns out to require a bit of subtlety, so let's consider a simpler question. Note that in the case $N = 1$, then the covariance matrix reduces to the scalar variance $\sigma^2$. We can then simplify the probability density function to

\[ p_1(x) = \frac{1}{\sigma\sqrt{2 \pi}} \exp \left ( -\frac{1}{2} \left ( \frac{x-\mu}{\sigma} \right )^2 \right ) \]

How can we sample from this simpler distribution? Let's assume that we have a method to sample from the uniform distribution $\mathcal{U}$ on the interval $[0, 1]$. Let's draw two samples $U_1$ and $U_2$ from $\mathcal{U}$. Then let 

\begin{eqnarray*}
Z_0 &= \sqrt{-2 \ln U_1} \cos(2\pi U_2) \\
Z_1 &= \sqrt{- 2 \ln U_1} \sin(2 \pi U_2) \\
\end{eqnarray*}

We claim that $Z_0$, and $Z_1$ are both samples from the univariate Guassian distribution with mean $\mu = 0$ and variance $\sigma^2 = 1$. This technique is commonly referred to as the Box-Muller transform. If you haven't seen this technique before, it may be somewhat mysterious to understand. Let's break down the derivation a bit. It might help if we solve these equations for $U_1, U_2$ (the derivation is left as an exercise).

\begin{align*}
U_1 &= \exp\left(-\frac{1}{2}(Z_0^2 + Z_1^2) \right) \\
U_2 &= \frac{1}{2\pi} \arctan\left (\frac{Z_1}{Z_0}\right) \\
\end{align*}

Let $f:\mathbb{R}^2 \to \mathbb{R}^2$ be defined as

\[ f(x_1, x_2) = \begin{bmatrix}
    \exp\left(-\frac{1}{2}(x_1^2 + x_2^2) \right) \\
    \frac{1}{2\pi} \arctan\left (\frac{x_2}{x_1}\right) \\
    \end{bmatrix} \]
    
such that $f(Z_0, Z_1) = U_1, U_2$. You can show via a simple calculation (left to the exercises) that 

\[ |\mathcal{J}(f)| = \frac{1}{2\pi}\exp \left ( -\frac{1}{2} (x_1^2 + x_1^2) \right ) = p_1(x_1) p_2(x_2) \]

That is, $|\mathcal{J}(f)|$ is the product of the probability density functions of two one dimensional Gaussians!

\section{Introduction to Normalizing Flows}

The core idea of a normalizing flow is to transform a simple probability distribution such as a Gaussian distribution into more sophisticated function which can model a real world data distribution through the use of a suitable transformation function $g$. Suppose that
\begin{align}
    z &\sim \mathcal{N}(\mu, \Sigma) \\
    x &= f(z)
\end{align}
Then the probability density of $x$ is given by
\begin{align}
    P(x) &= P_z(f(x))\left | \det \mathcal{J}(f) \right |
\end{align}
We constrain the function $g$ to be one-to-one and onto. %(\textcolor{red}{TODO: Define these in mathematical basics.}). 
A core idea of using a deep network to train a normalizing flow is that we can consider a chain of such transformations
\begin{align}
    z &\sim \mathcal{N}(\mu, \Sigma) \\
    x_1 &= f_1(z) \\
    x_2 &= f_2(x_1) \\
\end{align}
Then the final distribution is given by
\begin{align}
    P(x_2) &= P_z(f(x))\left | \det \mathcal{J}(f_2) \right | \left | \det \mathcal{J}(f_1) \right |
\end{align}
The loss for the normalizing flow is simply given by the log-likelihood
\begin{align}
    \mathcal{L} &= -\log P_z[f_n(\dotsc (f_1(x))] - \sum_i \log \left | \det \mathcal{J}(f_i) \right |
\end{align}
%\textcolor{red}{\url{https://dmol.pub/dl/flows.html}}
We can transform this description into Physika code as follows
\begin{lstlisting}[language=python]
class AutoregressiveNetwork(params: $\mathbb{N}$, hidden_units: [$\mathbb{N}$], activations: String)

class MaskedAutoregressiveFlow(net: Module)

class TransformedDistribution(zdist: $\mathbb{R}^N$, bijection: $\mathbb{R}^N \to \mathbb{R}^N$)

class NormalizingFlow(num_layers : $\mathbb{N}$, ndims: $\mathbb{N}$):
bijects = []
# loop over desired bijectors and put into list
for i in num_layers:
anet = AutoregressiveNetwork(
    params=ndim, hidden_units=[128, 128], activation="relu")
ab = MaskedAutoregressiveFlow(anet)
bijects.append(ab)
# Now permute 
permute = Permute([1, 0])
bijects.append(permute)
chained = Chain(bijects)

def $\lambda$(zdist):
# make transformed dist
td = TransformedDistribution(zdist, bijector=chained)
\end{lstlisting}

\begin{lstlisting}[language=python]
x : $\mathbb{R}^{N, 2}$
log_prob = log_prob(x)
model = Model(x, log_prob)
def neg_loglik(log_prob):
return -log_prob
\end{lstlisting}
% # now we prepare model for training model.compile(optimizer=Adam(1e-2), loss=neg_loglik)
\section*{Exercises}

\begin{enumerate}
    \item Let $U_1$ and $U_2$ be two samples from the uniform distribution on the interval $[0, 1]$. Let 
    \begin{eqnarray*}
Z_0 &= \sqrt{-2 \ln U_1} \cos(2\pi U_2) \\
Z_1 &= \sqrt{- 2 \ln U_1} \sin(2 \pi U_2) \\
\end{eqnarray*}
    Solve these equations for $U_1$ and $U_2$. Prove that 
    \begin{align*}
U_1 &= \exp\left(-\frac{1}{2}(Z_0^2 + Z_1^2) \right) \\
U_2 &= \frac{1}{2\pi} \arctan\left (\frac{Z_1}{Z_0}\right) \\
\end{align*}

\item Let $f:\mathbb{R}^2 \to \mathbb{R}^2$ be defined as

\[ f(x_1, x_2) = \begin{bmatrix}
    \exp\left(-\frac{1}{2}(x_1^2 + x_2^2) \right) \\
    \frac{1}{2\pi} \arctan\left (\frac{x_2}{x_1}\right) \\
    \end{bmatrix} \]
    
    Prove that $|\mathcal{J}(f)| = \frac{1}{2\pi}\exp(-\frac{1}{2} (x_1^2 + x_2^2))$
\end{enumerate}

%\section{Normalizing Flows}


\end{document}